% Figure: classical (thermoviscous) absorption coefficient of air versus
% frequency on log-log axes, showing the f^2 slope, with the
% characteristic absorption length 1/alpha on the right-hand reading.
\begin{figure}[htbp]
\centering
\begin{tikzpicture}
\begin{axis}[
  width=12cm, height=7.5cm,
  xlabel={frequency (Hz)},
  ylabel={classical absorption $\alpha_{\mathrm{cl}}$ (dB/km)},
  xmode=log, ymode=log,
  xmin=100, xmax=1000000,
  ymin=1e-3, ymax=1e5,
  axis lines=left,
  domain=100:1000000, samples=200,
  legend pos=north west,
  legend cell align=left,
  every axis x label/.style={at={(current axis.right of origin)}, anchor=west},
  every axis y label/.style={at={(current axis.above origin)}, anchor=south},
  grid=major, grid style={enggray!20},
  tick label style={font=\scriptsize},
  label style={font=\small},
  legend style={font=\scriptsize},
]
% Classical absorption ~ 1.6e-10 * f^2 in Np/m for air, converted to dB/km:
% alpha[dB/km] = 8.686 * 1000 * 1.6e-10 * f^2 = 1.39e-4 * f^2 (per km, in dB).
\addplot[engblue, very thick] {1.39e-4 * x^2};
\addlegendentry{classical $\alpha_{\mathrm{cl}}\propto f^2$}
% mark 1 dB/km reference level
\addplot[enggray, thin, dashed] coordinates {(100,1) (1000000,1)};
\node[enggray, font=\scriptsize, anchor=south east] at (axis cs:900000,1.3) {1 dB/km};
\end{axis}
\end{tikzpicture}
\caption[Classical thermoviscous absorption of air]{Classical
thermoviscous absorption coefficient of air versus frequency, drawn from
the Stokes-Kirchhoff sum of viscous and heat-conduction losses, which
grows as the square of frequency. The $f^2$ slope is a straight line of
gradient two on the log-log axes. A tone at $1\,\si{\kilo\hertz}$ loses a
fraction of a decibel per kilometre to the classical mechanism alone,
while a $1\,\si{\mega\hertz}$ ultrasonic beam loses of order a hundred
decibels per metre, which is why airborne ultrasound is a short-range
tool. The plotted classical curve is a lower bound on real atmospheric
loss: molecular relaxation of nitrogen and oxygen adds a humidity
dependent excess through the audible and low-ultrasonic bands that lifts
the total above this line, the mechanism the earlier attenuation figure
plots for two humidities. The reciprocal of the absorption coefficient is
the $e$-folding distance of the pressure amplitude, so the same curve read
as $1/\alpha$ gives the range over which a beam of a given frequency stays
usable.}
\label{fig:vol03ch07:classical-absorption}
\end{figure}
